This report is a long document, and it is easy to get lost in. The following text aims to help you navigate to the right sections depending on your need. Unfortunately, this project had generated a substantial amount of information, so some key considerations may have been missed or skipped. This report won't teach you everything I know, but \textbf{it will place you in my shoes and teach you why I made my decisions}.
\nl
\autoref{sec:intro} is your extended version of the project proposal. It explores a generic background and rationale for the project. Following this, the project scope is broken down and clearly explained, outlining exactly what is and isn't part of the project. Project milestones illustrate the core outcomes from the project and the responsibilities section shows how each milestone is divided between project members. Lastly, the project timeline presents the expected progression of the project.
\nl
Before diving into the other explanations, it's worth quickly understanding what I'll be referring to when I say v0, v0.1, and v1. It was intended that the project would include two versions, the first being a prototype to confirm functionality (v0), and a flight prototype (v1). Unfortunately, the test board needed to be redesigned, so a revised version of v0 was created, v0.1.
\nl
\autoref{sec:c2sv0} explores the first draft of the project (C2Sv0) that used the MIPI CSI-2 camera interface, which was eventually shelved. This section critically explores each of the core system components and provides thorough justications for each selection. These components mainly map to those used in C2Sv0, but their considerations and justifications are still relevant. You may find that this version is often referred to as the \textit{"breakout board"} or \textit{"P, S, A"} meaning (P)rocessor, (S)ensor, and (A)daptor. The schematic and PCB designs were split into separate sections. This version was deprioritised because the only chips that support the CSI-2 interface have a BGA pitch too fine for manufacturing within current budget limitations. Although this board was not procured, many circuits were recycled for the succeeding version (v0.1) and the content is still worth a read, especially for the CSI-2 interface information.
\nl
\autoref{sec:c2sv0.1} investigates the second draft of the project (C2Sv0.1) supported the DCMI camera interface. Unlike CSI-2, DCMI is older, so the chips that support the interface have much larger packages which were within our manufacturing capabilities. Like v0, the schematics and PCBs are split into different sections. Errors and fixes are mentioned as a remark for each section. \autoref{sec:manufacturing-v0.1} outlines the processes involved in manufacturing and developing the hardware. This involes where the components where sourced from, how the components were placed, as well as the 3D printed designs. \autoref{sec:v0.1-coding} dives into programming the hardware from the previous section. It focuses on coding each individual subsystem first.
\nl
\autoref{sec:c2sv1} follows the development of C2Sv1, which was the first version that demonstrated a complete image capture with the OV7670 camera module. This version explored a basic systems-engineering approach in \autoref{sec:v1-planning} to lay strong foundations and eliminate further errors throughout the hardware development process.
\autoref{sec:v1-schematic} and \autoref{sec:v1-pcb} explored the hardware aspects of the project.
\autoref{sec:v1-software} focused on the software development of the project, investigating real-time operating systems, memory controllers, and packet structure.
Lastly, \autoref{sec:v1-testing} detailed the results from testing, which included image quality and power consumption results. 
\nl
\autoref{sec:c2sv1.1} explores C2Sv1.1, the HAB- and FlatSat-compatible version. Unlike the previous versions, which used the STM32U5 series microcontroller, v1.1 housed a more powerful H7A3 chip to handle long compression times.
The development period of this version was very short to adhere with deadlines for HAB2, but still executed a more in-depth systems approach than v1 which served useful when making design decisions on short timelines in \autoref{sec:v1.1-planning}.
\autoref{sec:v1.1-schematic} and \autoref{sec:v1.1-pcb} encapsulated the hardware development of the version, including last-minute design changes.
The short timeline derailed the traditional manufacturing process of the board which was outlined in \autoref{sec:v1-manufacturing}. Moreover, the software implementation can be found in \autoref{sec:v1-software}, and testing and results found in \autoref{sec:v1-testing}.
\nl
The Appendices (\autoref{sec:appendices}) are supplementary materials that are non-essential to the project, but are a useful read. For example, some of the research is included to fill in gaps of knowledge. To the average reader, this is not-critical, but can be insightful.